\begin{abstract}
\noindent Answer engines such as ChatGPT and Perplexity answer questions by naming a few organisations and citing
sources. Answer engine optimisation (AEO) has focused on editing page content; we argue that access comes first, and
test both. Study~1 measures access barriers on \NSites{} organisational websites in \NSectors{} sectors and
\RandSites{} randomly sampled popular sites. One organisational site in three (\BarAny\%; \BarNoEdge\% without
refusals at the network edge) has a barrier to answer engines, usually an implicit one such as a wildcard robots.txt
rule, and sites refuse AI training crawlers far more often than answer agents. Study~2 is a pre-specified experiment with fictional brands: 14 page treatments in 12 hosted language models,
replicated on \FlNModelsTwo{} models behind widely used assistants. With the pages in
the model's context, a third party's recommendation raised the chance that a brand was named by \EMenAllThirdPartyU{}
percentage points, more than any edit to the brand's own page; citations and statistics added \EMenAllCitationsU{}
and \EMenAllStatisticsU{} points and were passed on to readers, repeated keywords also helped, and a hidden
instruction worked on most models but not on Anthropic's. When \AgNModelsW{} models had to search a pool of pages
themselves, the brand was named in \AgBaseMen\% of baseline answers instead of \AgBaseMenCtx\%, evidence edits and the
hidden instruction lost their effect because the page was seldom opened, and text matching the question made pages
far easier to find. In a live audit, ChatGPT's free tier cited no sources. We release the toolkit,
benchmark and data.
\end{abstract}
